\documentclass[11pt]{article}
\usepackage{icmlko}
\begin{document}

\title{CN 은 못 건넌 게 아니라 낮은 데 있었다 --- 여덟 번 만에,
질문을 뒤집으니 답이 나왔다}
\author{월드모델 실험실}
\date{노트 444 · 2026년 8월 2일}
\maketitle

\begin{abstract}
노트 440 이 연 물음이 여덟 번째다. CN 만화는 \texttt{entry\_friction} 이
\textbf{아예 없어} 사실상 이미 가려져 있는데도 $+0.2249$ 에 머물고, 같은
축을 가린 KR 은 $+0.5336$ 이다. 그 사이 $0.31$ 을 놓고 표본 크기 · 국가 ·
언어 · 연재 형식 · 축 개수 · 축 정체까지 일곱 가지가 죽었다.

이번엔 \textbf{질문을 뒤집었다}. 여태 물은 것은 전부 ``무엇이 전이를
막나''였다. 대신 물었다 --- \textbf{CN 자료 자체로 낼 수 있는 최대가
얼마인가.} 짝의 \emph{자기 행만으로} 교차검증 모형을 세워 천장을 쟀다.

\begin{center}
\begin{tabular}{lrrrr}
\hline
짝 & 축 & 천장 & 전이(가림) & 전이/천장 \\
\hline
KR 만화 ($n{=}1{,}716$) & 4 & $+0.7152$ & $+0.5336$ & $0.75$ \\
비게임 앱 ($n{=}1{,}600$) & 5 & $+0.6177$ & $+0.2635$ & $0.43$ \\
\textbf{CN 만화} ($n{=}352$) & 3 & $\mathbf{+0.3619}$ & $+0.2249$ & $\mathbf{0.62}$ \\
\hline
\end{tabular}
\end{center}

\textbf{CN 의 천장은 KR 의 절반이다.} 그리고 전이를 천장으로 나누면
CN 의 $0.62$ 는 KR 의 $0.75$ 와 앱의 $0.43$ \emph{사이}에 얌전히 앉는다 ---
\textbf{튀는 짝이 아니다.} CN 은 전이를 못 한 게 아니라 \textbf{갈 데까지
갔다}. 여덟 번째 설명이 처음으로 산다.
\end{abstract}

\begin{figure}[t]
\centering
\includegraphics[width=\textwidth]{figs/f1.pdf}
\caption{\textbf{가} --- 연한 막대가 짝이 자기 자료로 낼 수 있는 천장,
안쪽 진한 막대가 실제 전이. CN 은 막대 자체가 낮다. \textbf{나} ---
노트 443 의 몇 샷을 걸면 뒤집을 축이 있는 두 짝이 $0.75 \cdot 0.43$ 에서
$0.84 \cdot 0.78$ 로 \emph{모인다}. 가장 덜 찼던 앱이 가장 많이 오른다.}
\end{figure}

\section{천장이 낮은 이유는 태깅이다}

CN 은 축이 \textbf{셋}이다(\texttt{entry\_friction} 도 \texttt{media\_push}
도 없다). 남은 셋도 짝 안 부호가 KR 보다 일제히 약하다 ---
\texttt{target\_breadth} $+0.315$ 대 $+0.471$, \texttt{venue\_prominence}
$+0.045$ 대 $+0.162$, \texttt{goods\_scale} $+0.124$ 대 $+0.389$. 같은
\texttt{manga\_records} 수집인데 \textbf{CN 행에서 태거가 건진 신호가 적다}.

수집의 성질도, 나라도, 언어도 아니었다. \textbf{자료에 든 것이 적었을
뿐이다.}

\section{그래서 전이는 천장으로 나눠야 견줄 수 있다}

이 노트가 규약 하나를 만든다. \textbf{집 밖 짝끼리 견줄 때는 전이 값이
아니라 전이/천장 비로 견준다.} 짝마다 천장이 두 배씩 다른데 날것을 나란히
놓는 것은 \emph{키가 다른 사람의 점프 높이를 바닥에서 재는 것}이다.
노트 440 의 ``집 밖이 한 덩어리가 아니다''는 관측은 \textbf{반쯤 맞았다}
--- 날것으로는 $+0.047 \sim +0.53$ 로 흩어지지만 \textbf{비로 보면
$0.43 \sim 0.75$ 로 훨씬 좁다}.

\section{덤 --- 오늘의 다른 결과를 되맞힌다}

비는 예보도 한다. 앱은 세 짝 중 \textbf{가장 덜 찬} 짝이었다($0.43$) ---
그러니 고칠 여지가 제일 크다. 실제로 오늘 노트 441 의 가림이 산 값도
(KR $+0.3174$ 대 앱 $+0.4188$), 노트 443 의 몇 샷이 산 값도 앱이 더
컸다. 그리고 몇 샷($k{=}40$)을 걸면 둘이 \textbf{$0.84$ 와 $0.78$ 로
모인다}.

\textbf{같은 자리로 모인다는 것이 이 세 노트가 옳았다는 제일 좋은 증거다.}
서로 다른 두 수집이, 서로 다른 병을 앓다가, 같은 처방 뒤에 자기 천장의
같은 비율에 선다.

CN 은 못 모인다 --- \textbf{뒤집을 축이 없다}. CN 을 올리려면 모형이
아니라 \textbf{태거}를 고쳐야 한다.

\section*{남은 것}

\begin{itemize}
\item \textbf{CN 태거.} \texttt{entry\_friction} 과 \texttt{media\_push} 가
CN 행에서 왜 안 나오나. 이건 모형 갈래가 아니라 수집 갈래다.
\item \textbf{천장을 판에도 적용하기.} 판 $0.4766$ 도 도메인마다 천장이
다를 텐데 지금은 날것으로 가중평균한다. 시장팝업 $+0.145$ 는 정말 나쁜
건가, 천장이 낮은 건가.
\item \textbf{$0.43 \sim 0.75$ 도 아직 넓다.} 비를 더 좁히는 게 남았다.
\item LODO 씨앗 셋$\to$아홉 --- 여섯 노트째 미룸.
\end{itemize}

\end{document}
